\chapter*{前言}
\addcontentsline{toc}{chapter}{前言}
卡尔·弗里斯顿

主动推断是一种理解有感知能力之行为的方式。你正在阅读这些文字这一事实，本身就意味着你正在以一种特定方式进行主动推断，也就是在主动采样世界，因为你相信自己会从中学到些什么。你的眼睛之所以在页面上来回探查，只是因为这种行动能够消解你接下来会看到什么，以及这些词语究竟传达什么的不确定性。简言之，主动推断把“行动”放进了“知觉”之中，于是知觉被视为一种感知推断，或一种假设检验。主动推断更进一步，把规划也视为推断，也就是说，去推断你下一步会做什么，以便消解你对自身生活世界的不确定性。

为了说明主动推断的简洁性，以及我们试图解释的究竟是什么，请把指尖轻轻放在你的腿上，保持不动一两秒。现在，你的腿摸起来是粗糙还是光滑？如果你必须移动手指，才能唤起粗糙或光滑的感觉，那么你就已经发现了主动推断的一项基本原理。感受，就是触探；看见，就是去看；听见，就是去听。这种触探未必一定是外显的，我们也可以通过把注意力导向这里或那里来进行隐性的行动。简言之，我们并不只是试图理解自己的感觉；我们还必须主动创造自己的感官世界。在接下来的内容中，我们会看到为什么事情必须如此，也会看到为何我们所感知、所做、所规划的一切，都受一个存在性的根本命令所统摄，那就是“自我证成”。

主动推断并不只是关于阅读，或关于以认识为目的的信息觅取。从一种观点看，凡是生物乃至粒子，只要存在，就都在进行某种主动推断。这听起来也许是一个很强的主张；但它揭示了这样一个事实：主动推断承继自自由能原理，而这一原理将存在等同于自我证成，并将自我证成等同于一种具身生成的推断。不过，本书并不关心有感知系统的物理学本身。它关注的是，这种物理学对于理解大脑如何运作意味着什么。

这种理解绝非易事，几千年的自然哲学与几百年的神经科学已经充分说明了这一点。尽管我们可以在关于自组织行为的第一性原理解释中找到主动推断的根基，例如与哈密顿最小作用原理相近的变分原理，但当我们追问某一个具体大脑是如何工作的、又为何不同于另一个大脑时，第一性原理并不能帮上太多忙。比如，接受自然选择进化论，并不能丝毫帮助我们理解为什么我有两只眼睛，或者为什么我会说法语。本书的目标，是利用这些原理为神经科学与人工智能中的关键问题搭建脚手架。为此，我们必须超越原理本身，真正进入这些原理所作用其上的机制层面。

因此，主动推断以及与之相伴的贝叶斯机制，其作用是为我们如何感知、规划和行动的问题提供一个框架。关键在于，它并不试图取代行为心理学、决策理论、强化学习等其他框架。恰恰相反，它希望在一个统一框架中吸纳所有那些已经被证明卓有成效的方法。在下文中，我们将特别关注如何把心理学、认知神经科学、生成行动论、动物行为学等领域中的关键构念，与主动推断中的信念更新计算及其相关过程理论连接起来。

所谓过程理论，指的是解释信念更新如何在具身大脑乃至更广泛系统中的神经过程和其他生物物理过程中得以实现的理论。迄今为止，主动推断领域已经提供了一套相当直接的计算架构与仿真工具，既可用于建模一个运作中的大脑的不同方面，也能让研究者检验关于不同计算架构的假设。然而，这些工具只解决了问题的一半。主动推断的核心是一种生成模型，也就是对如下关系的概率表征：外部世界中那些不可观测的原因，如何生成可观测的结果，也就是我们的感觉。如何构造出恰当的生成模型，使之能够充分解释某个实验对象或某种生物的有感行为，这才是最大的挑战。

本书正是试图说明如何应对这一挑战。第一部分将搭建基本思想与形式化工具，第二部分则说明如何在实践中运用这些工具。简言之，本书面向那些希望使用主动推断来模拟和建模有感行为的人，无论他们的目的在于科学研究，还是可能在于人工智能。因此，本书聚焦于那些理解并实现一套主动推断方案所必需的思想与方法，而不会让读者被一方面有关有感系统物理学的问题、另一方面有关哲学的问题分散注意力。

\section*{卡尔·弗里斯顿的一则说明}

我得坦白一件事。这本书并不是我写的，至少大部分不是。更准确地说，别人不允许我多写。这本书的写作目标要求一种简洁、清晰的文风，而这并非我所擅长。尽管他们允许我偷偷塞进几个自己喜欢的词，但接下来的内容主要证明了托马斯和乔瓦尼对相关问题的深刻理解，以及同样重要的，他们在各种意义上的“心智理论”。

\section*{致谢}

我们衷心感谢朋友和同事们给予的宝贵帮助，尤其感谢伦敦大学学院威康信托人类神经影像中心理论神经生物学小组现在和过去的成员，感谢意大利国家研究委员会认知科学与技术研究所“行动中的认知”（Cognition in Action, CONAN）实验室的同仁，也感谢众多国际合作者；他们都对本书所呈现思想的发展发挥了关键作用。这个年轻却不断壮大的共同体在智识支持与持续激励方面都极为慷慨。与此同时，我们还诚挚感谢 MIT Press 的 Robert Prior 和 Anne-Marie Bono 在本书准备过程中给予我们的陪伴与指导，也感谢 Jakob Hohwy 以及其他审稿人提出的宝贵意见。最后，我们感谢为本研究提供资金支持的资助机构：KJF 获得了 Wellcome Trust 首席研究员基金资助（Ref: 088130/Z/09/Z）；GP 获得了欧洲研究委员会第 820213 号资助协议（ThinkAhead）以及欧盟“地平线 2020”研究与创新框架计划第 945539 号专项资助协议（Human Brain Project SGA3）的支持。
